\section{Preliminaries}\label{sec:prelim}

This section contains some background material referenced throughout the derivations throughout the thesis. By introducing the building blocks in this section, allowing the subsequent the presented arguments to remain focused on their primary intention. Having some known results as preliminaries will help us keep our derivations cleaner.

The central mathematical object in our derivations will be the (P)PGF of the model in question, as well as the idle probability of the system for some corollaries.


\subsection{\texorpdfstring{$M/M/1$}{M/M/1} queue}\label{ssec:mm1}

The model described here is the canonical $M/M/1$ queue, well studied in multiple textbooks and manuals and considered the absolute foundation of queueing theory. Here, we will use \cite{adan2002queueing} as our reference, but there is a lot of good quality literature and sources about it. 
Customers arrive at the system according to a Poisson process —i.e., the inter-arrival times are exponentially distributed with parameter $\lambda$—, service times are exponentially distributed with parameter $\mu$, and a single server operates in first-come-first-serve (FCFS) —also known as first-in-first-out (FIFO)— order. The system is stable if and only if the system's load $\rho = \lambda/\mu < 1$. 

Let $\pi_n$ be the limiting probability of finding $n$ customers in the system. Under stability, the detailed balance equations $\lambda\,\pi_n = \mu\,\pi_{n+1}$ ($n \geq 0$) together with the normalisation $\sum_{n=0}^{\infty}\pi_n=1$ give us the geometric steady-state distribution of the number of customers in the system:
\begin{equation}
    \pi_n = (1-\rho)\,\rho^n, \quad n \ge 0. \label{eq:mm1:steadystate}
\end{equation}
By the PASTA property (\S\ref{ssec:priority}), this also corresponds to the distribution seen by the arriving customers. Let $L$ be the random variable that tells us how many customers—including the one that might be in service—are in the system. The expected number of customers in the system is the following:
\begin{equation*}
    \mathbb{E}[L] = \frac{\rho}{1-\rho}.
\end{equation*}
From this, the expected sojourn time $W$—the time that elapses between a customer's arrival and their departure after service completion—is derived by Little's Law ($\mathbb{E}[L] = \lambda\mathbb{E}[W]$):
\begin{equation*}
    \mathbb{E}[W] = \frac{\mathbb{E}[L]}{\lambda} = \frac{1/\mu}{1-\rho}.
\end{equation*}

Another concept, more rare but highly relevant for the content that will be discussed in this thesis, is the Busy Period ($BP$) of the system, i.e., the time elapsed between the arrival of a customer to an idle system and the departure of a customer who leaves the system empty again after service completion. The Laplace-Stieltjes Transform (LST) and the expectation of this system's busy period are known:
\begin{align}
    \widetilde{BP}(s) &= \frac{(\mu+\lambda+s) - \sqrt{(\mu+\lambda+s)^2 - 4\mu\lambda}}{2\lambda},\label{eq:mm1:bp_lst} \\
    \mathbb{E}[BP]    &= \frac{1}{\mu(1-\rho)}. \label{eq:mm1:bp_mean}
\end{align}

% ─────────────────────────────────────────────────────────────────────────────
\subsection{\texorpdfstring{$M/M/1{+}M$}{M/M/1+M} queue (Erlang-A)}\label{ssec:erlangA}

Another relevant queueing model for this thesis is the so-called $M/M/1+M$ model, also known as the \emph{Erlang A} model. It's based on the $M/M/1$ queue —Poisson arrivals at a rate $\lambda$ and exponentially distributed waiting times with parameter $\mu$-, but adds customer impatience to it: each waiting customer can independently abandon at an exponentially distributed waiting time with parameter $\theta$. The key difference of this model with respect to the $M/M/1$ is that the $M/M/1+M$ does not require $\rho=\lambda / \mu \leq 1$ to be stable: it is positive recurrent for all $\lambda, \mu, \theta > 0$.

The steady state distribution follows, naturally, from its balance equations. At state $n \geq 0$ —for $n$ being the number of customers in the \emph{system}— we have
\begin{equation}
    \lambda\,\pi_n = (\mu + n\theta)\,\pi_{n+1}, \qquad n \ge 0.
    \label{eq:erlangA:balance}
\end{equation}
Iterating yields
\begin{equation}
    \pi_n = \pi_0\,\frac{(\lambda/\theta)^n}{(\mu/\theta)^{(n)}} \quad n \ge 0;
    \label{eq:erlangA:steadystate}
    \quad \text{where} \quad a^{(n)} = a(a+1)\cdots(a+n-1)
\end{equation}
Normalisation gives the non-trivial idle probability $\pi_0 = [{}_1F_1(1;\,\mu/\theta;\,\lambda/\theta)]^{-1}$ (see \S\ref{ssec:kummer} for the Kummer function ${}_1F_1$).

The \emph{Busy Period} of the $M/M/1+M$ defined here is far more complex than the one expressed in the $M/M/1$ queueing system. The LST and the expectation are given by
\begin{align}
    \begin{split}
    \widetilde{B}(s) = \cfrac{\mu}{\lambda_1+\mu+s
      -\cfrac{\lambda_1(\mu+\theta_1)}{\lambda_1+\mu+\theta_1+s
      -\cfrac{\lambda_1(\mu+2\theta_1)}{\lambda_1+\mu+2\theta_1+s-\dotsb}}},
    \label{eq:erlangA:bp_lst}
    \end{split} \\
    \begin{split}
        \mathbb{E}[B] = \dots
        \label{eq:erlangA:bp_exp}
    \end{split}
\end{align}
% TODO: Fill in the expected busy period
% TODO: Explain its integral representation and its Kummer representation

% ─────────────────────────────────────────────────────────────────────────────
\subsection{Pollaczek--Khinchine formula}\label{ssec:pk}

The $M/G/1$ queue —explained in \cite{adan2002queueing}— is a single-server queue with Poisson arrivals at a rate $\lambda$ but with generally distributed i.i.d.\ service times $B$ with mean $\mathbb{E}[B]$

The $M/G/1$ queue~\cite{adan2002queueing} is a single-server queue with Poisson arrivals
at rate $\lambda$ and generally distributed i.i.d.\ service times $B$ with mean and LST $\widetilde{B}(s) = \mathbb{E}[e^{-sB}]$. 
For stability purposes, the system load $\rho = \lambda\,\mathbb{E}[B] < 1$ is required.
Under these two conditions —Poisson arrivals and stability— the Pollaczek--Khinchine (PK) formula
gives the PGF of the stationary number of customers in the system:
\begin{equation}
    L(z) \;=\; \frac{(1-\rho_B)(1-z)\,\widetilde{B}(\lambda(1-z))}{\widetilde{B}(\lambda(1-z)) - z}.
    \label{eq:pk:pgf}
\end{equation}

% ─────────────────────────────────────────────────────────────────────────────
\subsection{Confluent hypergeometric function}\label{ssec:kummer}

The \emph{Kummer confluent hypergeometric function of the first kind} function will come up in some of the derivations of this thesis, and it reads as follows:
\begin{equation}
    {}_1F_1(a;\,b;\,z) \;=\;
        \sum_{n=0}^{\infty} \frac{a^{(n)}}{b^{(n)}\,n!}\,z^n,
    \qquad a^{(n)} = \prod_{k=0}^{n-1}(a+k),\quad a^{(0)}=1,
    \label{eq:kummer:series}
\end{equation}
where $b$ is not $0$ or a negative integer —i.e. $b \notin \{0,-1,-2,\dots\}$— and the series converges for all
$z\in\mathbb{C}$.
The function ${}_1F_1(a;\,b;\,z)$ is the unique solution of \emph{Kummer's equation}~\cite{dlmf}:
\begin{equation}
    z\,u'' + (b - z)\,u' - a\,u = 0, \label{eq:kummer:ode}
\end{equation}
normalized by ${}_1F_1(a;\,b;\,0) = 1$. The \emph{Kummer series} (Equation \eqref{eq:kummer:series}) converges
for all $z\in\mathbb{C}$ and $b \notin \{0,-1,-2,\ldots\}$.

For this thesis, two identities related to it are used.
First, the \emph{Euler integral representation}, for $a,b>0$ we have:
\begin{equation}
    {}_1F_1(a;\,a+b;\,z) \;=\; \frac{1}{B(a,b)} \int_0^1 t^{a-1}(1-t)^{b-1} e^{zt}\,dt, 
    \label{eq:kummer:euler}
\end{equation}
where $B$ is the \emph{Beta distribution} $B(a,b)=\Gamma(a)\Gamma(b)/\Gamma(a+b)$, and $\Gamma$ is the \emph{Gamma function}.
Secondly, the \emph{Contiguous-function ratio}:
\begin{equation}
    \frac{{}_1F_1(a+1;\,b;\,z)}{{}_1F_1(a;\,b;\,z)}.
    \label{eq:kummer:ratio}
\end{equation}


% ─────────────────────────────────────────────────────────────────────────────
\subsection{First-order linear ODEs and PDEs}\label{ssec:ode}

Since jockeying and abandoning add state-dependent rates to our systems, we will have to deal with derivative terms to solve for our PGF's. Here follow some introductions to the methods employed to solve diverse typeS of ODE's and PDE's that will play a role in this thesis.

\subsubsection{Integrating factor (variation of parameters)}

This method is standard to tackle —for each fixed $y—$ a \emph{first-order linear ODE in $x$}, and leads to a closed-form solution, in our case, for $P(x,y)$. Such an equation has standard form
\begin{equation}
    \frac{dP}{dx} + p(x;\,y)\,P \;=\; q(x;\,y). \label{eq:ode:standard}
\end{equation}
The integratimg factor $\mu_I(x;y)$, defined by $\mu_I(x;y) = \exp\bigl(\int_0^x p(t;y)\,dt\bigr)$, converts the standard form (Equation \eqref{eq:ode:standard}) to $\frac{d}{dx}[P\cdot\mu_I] = q\cdot\mu_I$, yielding
\begin{equation}
    P(x,y) \;=\; \frac{1}{\mu_I(x;\,y)}
        \left[C(y) + \int_0^x q(t;\,y)\,\mu_I(t;\,y)\,dt\right].
    \label{eq:ode:variation}
\end{equation}
For $P$ bounded, $\mu_I(0;y)=0$, we have $C(y)=0$, giving
\begin{equation}
    P(x,y) \;=\; \frac{1}{\mu_I(x;\,y)}
        \int_0^x q(t;\,y)\,\mu_I(t;\,y)\,dt.
    \label{eq:ode:variation}
\end{equation}


\subsubsection{Method of characteristics}\label{sssec:moc}

This method is standard for solving a \emph{first-order linear PDE in both $x$ and $y$}, simultaneously. The Method of Characteristics reduces such an equation to a family of ODE's along the so-called \emph{characteistic curve} in the $(x,y)$-plane, and each of those collapse to an ODE in a single parameter, namely $t$. The general solution, thus, is obtained by parameterising one variable in terms of the other and integrating along the characteristics. this, together with an arbitrary function determined by the boundary data, yields a general solution.

A PDE of the form
\begin{equation*}
    a \frac{\partial P}{\partial x} + b \frac{\partial P}{\partial y} = cP + d
\end{equation*}
is solved by tracing the \emph{characteristic curve} —$(x(t),y(t))$, satisfying
\begin{equation}
    \frac{dx}{dt} = a, \quad \frac{dy}{dt} = b.
\end{equation}
Along each of these curves, the function $P(x(t), y(t))$ satisfies:
\begin{equation}
    \frac{dP}{dt} = cP + d
\end{equation}
which is an ODE in the single variable $t$.

